\documentclass[11pt,a4paper]{article}
\usepackage[utf8]{inputenc}
\usepackage[spanish]{babel}
\usepackage[margin=1in]{geometry}
\usepackage{enumitem}
\usepackage{hyperref}
\usepackage{graphicx}
\hypersetup{
    colorlinks=true,
    linkcolor=blue,
    filecolor=magenta,
    urlcolor=cyan,
}

\newcommand{\comment}[1]{}

\usepackage{pdfpages}
% Configure page margins with geometry
\geometry{left=1.4cm, top=.8cm, right=1.4cm, bottom=1.8cm, footskip=.5cm}

\usepackage[usenames,dvipsnames]{xcolor}
\usepackage{tikz} 
\usetikzlibrary{calc, arrows.meta, intersections, patterns, positioning, shapes.misc, fadings, through,decorations.pathreplacing}

\definecolor{ColorOne}{named}{MidnightBlue}
\definecolor{ColorTwo}{named}{Dandelion}
\definecolor{ColorThree}{named}{Plum}




% Color for highlights
% russell Colors: russell-emerald, russell-skyblue, russell-red, russell-pink, russell-orange
%                 russell-nephritis, russell-concrete, russell-darknight, russell-purple
\colorlet{russell}{russell-black}

\setlist[itemize]{noitemsep, topsep=0pt}

\begin{document}

\section{Información Personal}
\begin{itemize}
    \item \textbf{Nombre:} Dr. Francisco Richter Mendoza
    \item \textbf{Nacionalidad:} Chilena
    \item \textbf{Dirección Académica:} Via la Santa 1, CH-6962 Viganello, Suiza
    \item \textbf{Teléfono:} +41 782419782
    \item \textbf{Email:} \href{mailto:richtf@usi.ch}{richtf@usi.ch}
    \end{itemize}
\end{itemize}

\section{Estudios}
%\subsection*{2.1 Estudios Superiores Universitarios Formales}
\begin{itemize}
    \item \textbf{Doctorado en Estadística Computacional}
    \begin{itemize}
        \item Institución: University of Groningen, Países Bajos
        \item Tesis: \textit{Inferring the Drivers of Species Diversification}
        \item Fechas: 2016--2021
    \end{itemize}
    \item \textbf{Ingeniería Civil Matemática}
    \begin{itemize}
        \item Institución: Universidad Técnica Federico Santa María, Chile
        \item Tesis: \textit{Mathematical Modelling Applied to Photometric Correction in Astronomical Observations}
        \item Fechas: 2006--2013
    \end{itemize}
\end{itemize}

%\subsection*{2.2 Actividades de Perfeccionamiento}
%\begin{itemize}
%    \item Taller de Redacción Científica, Università della Svizzera italiana, 2022.
%    \item Formación en Métodos Estadísticos Avanzados, Universidad de Chile, 2020.
%\end{itemize}



%\subsection*{3.2 Premios y Otras Distinciones}
%\begin{itemize}
%    \item Lista de Honor, Universidad Técnica Federico Santa María, [Año(s)].
%    \item Associate Editor, \textit{Statistica Neerlandica}, Wiley Online Library (2023--Presente).
%\end{itemize}



\section{Cronología Laboral}

\begin{figure}[h]
\centering
\input{timeline}
    \centering
    \includegraphics[width=0.5\linewidth]{}
    \caption{Cronología laboral. Click en texto para dirigirse a thesis, paginas, etc.}
    \label{fig:enter-label}
\end{figure}

Mi trayectoria profesional se ha enfocado en resolver problemas complejos mediante herramientas matemáticas y estadísticas avanzadas, integrando investigación aplicada y docencia:

\begin{itemize}
    \item \textbf{Trabajo Postdoctoral en Università della Svizzera italiana (2021--Presente):}  
    Investigador postdoctoral. Mis actividades incluyen:
    \begin{itemize}
        \item Desarrollo de modelos estadísticos para redes complejas y sistemas dinámicos.
        \item Investigación sobre diversificación dependiente de la diversidad y análisis filogenético.
        \item Impartición de cursos (ver mas abajo).
    \end{itemize}
    
    \item \textbf{Doctorado en la Universidad de Groningen (2016--2021):}  
    Supervisado por el Dr. Rampal Etienne y el Prof. Ernst C. Wit, mi investigación exploró la diversificación de especies mediante marcos matemáticos y estadísticos, integrando conceptos de sistemas dinámicos y biología evolutiva.

    \item \textbf{Evalueserve, Chile (2013--2016):}  
    Analista de riesgos operativos y científico de datos. Aplicación de teoría de valores extremos y modelamiento estocástico en la evaluación de riesgos financieros.

    \item \textbf{European Southern Observatory (ESO), Chile (2013):}  
    Asistente de investigación. Desarrollé mi tesis de Ingeniería Matemática bajo la supervisión del Dr. Fernando Selman, enfocándome en métodos para correcciones fotométricas en observaciones astronómicas.

 \end{itemize}
    










\section{Actividades Académicas}

\subsection{Docencia Lectiva}

\begin{itemize}
    \item \textbf{Università della Svizzera italiana (USI):}
    \begin{itemize}
        \item \textit{Analysis of Social Networks} (Master in Computational Science, Elective, 2nd year), Autumn 2024.
        \begin{itemize}
            \item Este curso explora cómo las conexiones entre individuos y grupos influyen en sus acciones, cubriendo principios fundamentales, métodos estadísticos y aplicaciones prácticas.
        \end{itemize}
        \item \textit{Introduction to Ordinary Differential Equations} (Master in Artificial Intelligence and Computational Science, Elective, 1st and 2nd year), Autumn 2024.
        \begin{itemize}
            \item Curso centrado en la importancia de las ecuaciones diferenciales ordinarias en el modelado de sistemas reales. Incluye métodos analíticos, técnicas numéricas y aplicaciones prácticas.
        \end{itemize}
        \item \textit{Probability \& Statistics} (Bachelor in Informatics, 2nd year), Autumn 2023 y 2024.
        \begin{itemize}
            \item Proporciona una base comprensiva en probabilidad y estadística, enfatizando su aplicación a fenómenos del mundo real mediante modelos estadísticos.
        \end{itemize}
        \item \textit{Stochastic Methods} (Master in Artificial Intelligence and Computational Science), Spring 2024.
        \begin{itemize}
            \item Curso avanzado en procesos estocásticos, combinando teoría y aplicaciones prácticas.
        \end{itemize}
    \end{itemize}
\end{itemize}

Para más detalles, se puede visitar la página: \href{https://franciscorichter.github.io/teaching/}{https://franciscorichter.github.io/teaching/}.

\subsection{Docencia Personalizada}
\begin{itemize}
    \item Supervisión de tesis de maestría:
    \begin{itemize}
        \item Mattia Colbertaldo - \textit{Machine Learning for Species Diversification Dynamics}, Politecnico di Milano, 2023. Fecha de defensa: 13 de junio, 2024. 
        \item Lodovico Mazzei - \textit{Stochastic Learning for Species Diversification Dynamics}, Università della Svizzera italiana, 2023. Fecha de defensa: 5 de septiembre, 2023. 
        \item Oscar Arturo Silva Castellanos - \textit{Inferring Phylogenetic Diversification Parameters Using Deep Learning}, Università della Svizzera italiana, 2023. Fecha de defensa: 27 de noviembre, 2023.
    \end{itemize}
\end{itemize}




\end{itemize}

\subsection{Otras Actividades Académicas}

\begin{itemize}
  \item  Associate Editor, \textit{\href{https://onlinelibrary.wiley.com/page/journal/14679574/homepage/editorialboard.html}{\textit{Statistica Neerlandica}}}, Wiley Online Library (2023--Presente).
  
\end{itemize}



\section{Investigación y Comunicaciones}

\subsection{Publicaciones}
\begin{itemize}
    \item Zhou Y, Abrahams S, Mendoza F, Donà M, Verbiest R. \textit{Stereo 3D cloud motion from tandem satellites}. Mathematics in Industry Reports, 2024. DOI: \href{https://doi.org/10.33774/miir-2024-51r3q}{10.33774/miir-2024-51r3q}.
    \item Vinciotti V, Wit EC, Richter F. \textit{Random graphical model of microbiome interactions in related environments}. Journal of Agricultural, Biological and Environmental Statistics, 2024.
    \item Hendriks KP, Bisschop K, Kortenbosch HH, Kavanagh JC, Larue AE, et al., Richter FJ. \textit{Microbiome and environment explain the absence of correlations between consumers and their diet in Bornean microsnails}. Ecology, 2021. DOI: \href{https://doi.org/10.1002/ecy.3237}{10.1002/ecy.3237}.
    \item Hendriks KP, Richter FJ, et al. \textit{Plant diets of land snail community members are similar in composition but differ in richness}. Journal of Molluscan Studies, 2021.
    \item Richter F, Janzen T, Hildenbrandt H, Wit EC, Etienne RS. \textit{Detecting phylodiversity-dependent diversification with a general phylogenetic inference framework}. bioRxiv, 2021.
    \item Richter F, Haegeman B, Wit EC, Etienne RS. \textit{Introducing a general class of species diversification models for phylogenetic trees}. Statistica Neerlandica, 2020.
    \item Hendriks KP, Bisschop K, Kavanagh JC, Kortenbosch HH, Larue AE, Richter FJ. \textit{Fieldwork to sample microsnails for diet and microbiome studies along the Kinabatangan River, Sabah, Malaysian Borneo}. Malacologist, 2019.
\end{itemize}

\subsection{Ponencias a Congresos Nacionales e Internacionales}
\begin{itemize}
    \item \textbf{2023:} International Conference on Complex Networks (CompleNet). Aveiro, Portugal. \textit{Detecting Complex Diversity-Dependent Diversification with A Novel Phylogenetic Inference Framework}.
    \item \textbf{2023:} World Conference on Natural Resource Modeling. Amsterdam, Países Bajos. \textit{Phylodiversity-Dependent Diversification: A General Framework for Investigating the Role of Phylogenetic Diversity}.
    \item \textbf{2022:} Conference on Complex Systems. Palma de Mallorca, España. \textit{Inferring the Drivers of Species Diversification Processes}.
    \item \textbf{2020:} World Conference on Natural Resource Modeling. Valparaíso, Chile. \textit{Including Phylodiversity in Diversity-Dependent Diversification Models}.
    \item \textbf{2018:} Mathematics for Planet Earth Meeting. Utrecht, Países Bajos. \textit{A Statistical Approach to Species Diversification Dynamics}.
    \item \textbf{2018:} Conference on Complex Systems. Thessaloniki, Grecia. \textit{Generalizing Species Diversification Models}.
    \item \textbf{2017:} Mathematical Models in Ecology and Evolution. Londres, Reino Unido. \textit{Generalizing Species Diversification Models}.
    \item \textbf{2017:} International Workshop on Statistical Modelling. Groningen, Países Bajos. \textit{A General Statistical Framework to Study the Diversification of Species}.
\end{itemize}


%\section{Actividades Profesionales}

\subsection{Proyectos}
A lo largo de mi carrera, he participado en diversos proyectos financiados que abarcan modelamiento matemático, estadística y ciencias interdisciplinarias. A continuación, destaco algunos proyectos relevantes:

\begin{itemize}
    \item \textbf{STSM Grant (COST Action CA15109, 2020):}  
    Financiado por la European Cooperation in Science and Technology. Investigador principal en modelamiento estocástico aplicado a redes complejas. Colaboración entre la Universidad de Groningen y la Università della Svizzera italiana.
    
    \item \textbf{Mathematics for planet earth Grant (2016--2021):}  
    Proyecto de investigación financiado por la University of Groningen, enfocado en procesos de diversificación de especies y filogenética. Supervisores: Dr. Rampal Etienne y Prof. Ernst C. Wit.
    
    \item \textbf{Sparce Inference on Complex Networks (2021--Presente):}  
    Proyecto de modelamiento estadístico y análisis de redes complejas, financiado por la Università della Svizzera italiana. Investigador asociado al laboratorio del Prof. Ernst C. Wit.
\end{itemize}




\subsection{Investigación y Desarrollo}
\begin{itemize}
    \item Desarrollo del paquete \href{https://cran.r-project.org/web/packages/rgm/index.html}{\texttt{rgm} (Random Graphical Models)} para analizar redes biológicas y sociales.
    \item Colaboración con OASI en proyectos de modelado ambiental en tiempo real utilizando aplicaciones interactivas como \href{https://franciscorichter.shinyapps.io/EPIMOS/}{EPIMOS}.
\end{itemize}

\subsection{Divulgación}

\subsubsection{Workshops Organizados}

\begin{itemize}
    \item \textbf{2023: } \href{https://www.ci.inf.usi.ch/workshop-sparse-inference-on-complex-networks}{Sparse Inference on Complex Networks Workshop}. Lugano, Suiza.
    \item \textbf{2024: } \href{https://www.ci.inf.usi.ch/innodyn/}{Innovation in Dynamic Network Modelling Workshop}. Lugano, Suiza. 
\end{itemize}

\subsubsection{Divulgación audiovisual}

\begin{itemize}
    \item Mi tesis de doctorado en 5 minutos: \textit{How did species evolve?} Disponible en \href{https://www.youtube.com/watch?v=0RwCzOyzBXY}{YouTube}.
\end{itemize}

\section{Referencias}

Los siguientes profesores están disponibles para proporcionar cartas de recomendación bajo solicitud:

\begin{itemize}
    \item \textbf{Prof. Ernst C. Wit} \\
    Catedrático de Estadística y Ciencia de Datos \\
    \textit{Università della Svizzera italiana (USI), Suiza} \\
    Email: \href{mailto:ernst.jan.camiel.wit@usi.ch}{ernst.jan.camiel.wit@usi.ch}

    \item \textbf{Prof. Pedro Gajardo} \\
    Profesor Titular, Departamento de Matemáticas \\
    \textit{Universidad Técnica Federico Santa María, Chile} \\
    Email: \href{mailto:pedro.gajardo@usm.cl}{pedro.gajardo@usm.cl}

 

    \item \textbf{Prof. Berta Gabriela Olave Pereira} \\
    Docente de Matemáticas \\
    \textit{Instituto Nacional, Chile} \\
    Email: \href{mailto:b.olave.mat@institutonacional.cl}{b.olave.mat@institutonacional.cl}

\end{itemize}

%Las cartas de recomendación serán enviadas directamente por los evaluadores bajo solicitud.

\comment{

\section*{10. Antecedentes Adicionales}

A lo largo de mi carrera profesional y académica, he trabajado en diversos contextos donde el modelamiento matemático ha sido el núcleo de mi contribución. Este enfoque ha permitido abordar problemas complejos en astronomía, riesgos operativos, biodiversidad y estadística de redes.

\begin{itemize}
    \item \textbf{Universidad Técnica Federico Santa María (UTFSM), Chile:} 
    Colaboré con el \href{http://pgajardo.mat.utfsm.cl/}{Prof. Pedro Gajardo} en proyectos enfocados en el modelamiento matemático y optimización. Además, desarrollé mi tesis de pregrado bajo la supervisión del astrónomo de ESO, \href{https://www.sc.eso.org/~fselman/}{Dr. Fernando Selman}, trabajando en métodos matemáticos para correcciones fotométricas en observaciones astronómicas, mejorando datos obtenidos en telescopios de gran escala.
    
    \item \textbf{Evalueserve, Chile:} 
    Trabajé como analista de riesgos operativos en colaboración con \href{https://www.kansascityfed.org/research-staff/padma-sharma/}{Padma Sharma} y \href{https://business.ku.edu/people/trambak-banerjee}{Trambak Banerjee}. Aplicamos técnicas de teoría de valores extremos para el análisis de datos y la evaluación de riesgos en contextos financieros, diseñando herramientas para mitigar la incertidumbre en procesos empresariales.

    \item \textbf{Doctorado en la Universidad de Groningen, Países Bajos:} 
    Durante mi doctorado, trabajé bajo la supervisión del \href{https://www.rug.nl/staff/r.s.etienne/}{Dr. Rampal Etienne} y el \href{https://www.rug.nl/staff/e.c.wit/?lang=en}{Prof. Ernst C. Wit} en el desarrollo de modelos matemáticos para estudiar la diversificación de especies. Este trabajo incluyó la creación de marcos estadísticos para analizar procesos de diversificación dependiente de diversos tipos de diversidad, combinando ideas de sistemas dinámicos y filogenética.

    \item \textbf{Postdoctorado en Università della Svizzera italiana (USI), Suiza:} 
    Desde 2021, colaboro con el \href{https://www.ci.inf.usi.ch/research/statslab/people/}{Prof. Ernst C. Wit}, quien ha sido mi mentor durante los últimos 10 años. Nuestro trabajo se ha centrado en el modelamiento matemático y estadístico de sistemas complejos, particularmente en redes estadísticas y sus aplicaciones en biología, finanzas y ciencia de datos.
\end{itemize}



\textbf{Proyectos destacados:}  

        
\begin{itemize}
    \item \textbf{Modelos gráficos para interacciones en microbiomas:} Desarrollo de modelos gráficos aleatorios para analizar interacciones microbianas en ambientes relacionados, publicados en el \textit{Journal of Agricultural, Biological and Environmental Statistics} (2024).
    
    \item \textbf{Diversificación filogenética dependiente de la diversidad:} Introducción de un marco general de inferencia filogenética para detectar procesos de diversificación dependiente de la filodiversidad, presentado en \textit{bioRxiv} (2021).
    
    \item \textbf{Estudio de comunidades de caracoles terrestres en Borneo:} Análisis de la dieta y el microbioma en comunidades de caracoles terrestres para identificar correlaciones con su entorno. Este trabajo ha sido publicado en revistas como \textit{Ecology} (2021) y el \textit{Journal of Molluscan Studies} (2021).
    
    \item \textbf{Diversificación de especies:} Propuesta de una clase general de modelos de diversificación para árboles filogenéticos, publicada en \textit{Statistica Neerlandica} (2020).
    
    \item \textbf{Trabajo de campo en Malasia:} Estudios de campo realizados en el río Kinabatangan para muestrear microsnails y analizar su dieta y microbioma, con resultados publicados en el \textit{Malacologist} (2019).
    
    \item \textbf{Movimiento de nubes en 3D:} Colaboración internacional para estimar el movimiento tridimensional de nubes mediante satélites en tándem, con resultados publicados en \textit{Mathematics in Industry Reports} (2024).


   
\end{itemize}



}



\end{document}
